\textbf{Affirmation 1 : Fausse.}

	Le budget total pour un vélo est égal à $300+45+60 = 405$~\euro (document n°1), soit plus de 400~\euro.

	Donc, pour cinq vélos, le montant serait \np[\euro]{2025}, supérieur à \np{2000}~\euro, qui est le montant maximum prévu par la proviseure (document n°4).

	On ne pourra pas installer cinq vélos avec ce budget, l'affirmation est fausse.

\vspace{5mm}
\textbf{Affirmation 2 : Vraie.}

	En pédalant pendant une heure (60 minutes), on récupère \np[~kJ]{100} (document n°1).

	Il y a une relation de proportionnalité entre la durée du pédalage et l'énergie récupérée.

	Pour charger un téléphone portable, il faut \np[~kJ]{40} (document n°2):

	\begin{tabular}{c|c|c}
		durée&60  &$x$ \\\hline
		énergie&100  & 40 \\
	\end{tabular}\quad
	Il faut $\dfrac{40\times60}{100} = \dfrac{\np{2400}}{100} = 24$ minutes de pédalage pour récupérer \np[~kJ]{40}, et ainsi pouvoir charger un téléphone portable, donc l'affirmation est vraie.

\vspace{5mm}
\textbf{Affirmation 3 : Vraie.}

	Comme on choisit un élève au hasard, on est en situation d'équiprobabilité, donc les proportions sont assimilables à des probabilités.

	Parmi les 560 élèves favorables au projet, 180 sont des élèves de première (document n°3), donc la probabilité recherchée est $\dfrac{180}{560} = \dfrac{18}{56} = \dfrac{9}{28}$.

	On se souvient que $7\times4 = 28$ : le quart de 28 est donc 7, et donc $\dfrac{9}{28}$ est plus grand qu'un quart, l'affirmation est vraie.

\vspace{5mm}
\textbf{Affirmation 4 :}

	Deux conditions doivent être remplies pour que le projet soit lancé (document n°4):
	\begin{itemize}
		\item plus de la moitié des élèves du lycée doivent se déclarer favorables à cette initiative : c'est le cas, car le nombre d'élèves favorables est 560, ce qui est plus de la moitié des \np{1000} élèves du lycée (document n°3);

		\item au moins 60~\% des élèves de première doivent également y être favorables : c'est aussi le cas puisqu'il y a 300 élèves de première, dont 180 soit exactement 60~\% sont favorables au projet (document n°3), en effet, 60\,\% de 100 élèves, c'est 60 élèves, donc 60\,\% de 300 élèves, c'est le triple, soit $60 \times 3 = 180$.
	\end{itemize}
	Les deux conditions sont remplies, le projet pourra être lancé et l'affirmation est vraie.

\needspace{10\baselineskip}
